% Anhang --- zählt nicht zu den 10 bis 15 Seiten.
% Alles, was Beleg ist und nicht Argument, gehört hierher.

\appendix
\section{Anhang}

\subsection{Faktkatalog}
\label{app:faktkatalog}
\todo{Tabelle der untersuchten Sachverhaltstypen: Typ, Beispiel, erwartete
Änderungsrate, Quellen. Aus dem Team-Repo übernehmen, sobald es steht.}

\subsection{Extraktionsschema}
\label{app:schema}

Das Schema ist die technische Form des Sachverhalts-Tupels aus
Abschnitt~\ref{sec:methode}. Es hat drei Ebenen: das Dokument, den Fakt und die Fundstelle.
Tabelle~\ref{tab:schema} nennt die Felder.

\begin{table}[ht]
  \centering
  \caption{Felder des Extraktionsschemas}
  \label{tab:schema}
  \small
  \begin{tabular}{@{}>{\raggedright\arraybackslash}p{0.24\textwidth}>{\raggedright\arraybackslash}p{0.22\textwidth}c>{\raggedright\arraybackslash}p{0.34\textwidth}@{}}
    \toprule
    Feld & Typ & Pflicht & Regel \\
    \midrule
    doc\_id & Zeichenkette & ja & Kennung aus dem Datenblatt, z.\,B.\ \texttt{ord\_3ae} \\
    doc\_hash & SHA-256 & ja & bindet den Fakt an die eingefrorene Fassung \\
    entitaet & Zeichenkette & ja & kanonisch nach Aliasliste \\
    attribut & kontrolliertes Vokabular & ja & z.\,B.\ \texttt{bewerbungsfrist\_wise} \\
    wert\_raw & Zeichenkette & ja & wörtlich, wie im Dokument \\
    wert\_norm & Datum, Zahl, Text & nein & Ergebnis der Normalisierung \\
    einheit & Zeichenkette & nein & LP, Minuten, EUR \\
    scope & Zeichenkette & ja & Zielgruppe oder Programm; Vorgabe \texttt{alle} \\
    gueltig\_von & Datum & nein & Beginn des Gültigkeitsintervalls \\
    gueltig\_bis & Datum & nein & Ende des Gültigkeitsintervalls \\
    fundstelle\_text & Zeichenkette & ja & wörtliches Zitat \\
    fundstelle\_start & ganze Zahl & ja & Zeichenoffset im segmentierten Text \\
    fundstelle\_ende & ganze Zahl & ja & Zeichenoffset im segmentierten Text \\
    status & Aufzählung & ja & \texttt{ok} oder \texttt{verworfen} \\
    prompt\_version & Zeichenkette & ja & Reproduzierbarkeit des Laufs \\
    modell & Zeichenkette & ja & Name und Version des Modells \\
    \bottomrule
  \end{tabular}
\end{table}

Die Pflichtfelder sind so geschnitten, dass jeder Fakt ohne Rückfrage prüfbar ist:
\texttt{doc\_hash} und die beiden Offsets verweisen auf eine Stelle in einer eingefrorenen
Datei, nicht auf ein Dokument, das sich inzwischen geändert haben kann. \texttt{wert\_norm}
ist bewusst nicht Pflicht. Die Normalisierung darf scheitern, ohne den Fakt zu verlieren:
Der Rohwert bleibt erhalten, der Fall wird als nicht normalisierbar markiert und in der
Stichprobe manuell geprüft. Ein Fakt ohne Fundstelle wird dagegen nicht geschrieben --- er
gilt als nicht extrahiert.

Die Prüfung erfolgt beim Schreiben. Ein Datensatz, der ein Pflichtfeld nicht erfüllt, wird
verworfen und protokolliert, nicht korrigiert; damit ist die Datenbank eine Teilmenge des
tatsächlich Gelesenen und keine Interpretation. \texttt{wert\_raw} und \texttt{wert\_norm}
stehen nebeneinander, sodass sich jede Regel aus Abschnitt~\ref{sec:methode} am Einzelfall
nachvollziehen lässt; wird eine Regel geändert, entstehen neue normalisierte Werte, ohne
dass die Extraktion wiederholt werden muss. Der Hash schließlich trennt Dokumentversionen:
Ein zweiter Abruf derselben Adresse ist eine neue Fassung und keine Korrektur der alten ---
genau die Unterscheidung, die für Versionsartefakte gebraucht wird.

\subsection{Verwendete Sprachmodelle und Prompts}
\label{app:modelle}
\todo{Modellname, Version, Temperatur, Prompt-Version, Datum. Pflicht, wenn ein
Sprachmodell im Spiel ist --- sonst ist die Arbeit nicht reproduzierbar.}

\subsection{Fehlertaxonomie}
\label{app:fehlertaxonomie}

Die Taxonomie beschreibt Fehler des Verfahrens, nicht Abweichungen der Quellen. Sie ist von
der Typologie in Abschnitt~\ref{sec:typologie} zu unterscheiden: Dort geht es darum, was ein
Fund bedeutet; hier darum, wodurch das Verfahren danebenliegt. Tabelle~\ref{tab:fehler}
nennt die fünf Klassen.

\begin{table}[ht]
  \centering
  \caption{Fehlerklassen des Verfahrens}
  \label{tab:fehler}
  \small
  \begin{tabular}{@{}>{\raggedright\arraybackslash}p{0.20\textwidth}>{\raggedright\arraybackslash}p{0.44\textwidth}>{\raggedright\arraybackslash}p{0.26\textwidth}@{}}
    \toprule
    Klasse & Definition & Folge für die Kennzahl \\
    \midrule
    Extraktion & Der Wert steht so nicht im Dokument; das Zitat trägt ihn nicht & falscher Widerspruch (Precision) oder fehlender Fakt (Recall) \\
    Erfundene Angabe & Wert ohne Deckung im Text; trotz Zitatpflicht möglich, wenn das Zitat nicht zum Wert passt & Precision \\
    Normalisierung & richtig gelesen, falsch vereinheitlicht, etwa 1,5 Stunden als 90 statt 150 Minuten & Precision (falscher Widerspruch) \\
    Gruppierung & verschiedene Sachverhalte zusammengeworfen oder derselbe getrennt & Precision und Recall \\
    Scope & Zielgruppe übersehen, etwa EU-Bewerbende gegen alle & Precision \\
    Urteil & Widerspruch gesehen, wo keiner ist, oder einen übersehen & Precision bzw.\ Recall \\
    \bottomrule
  \end{tabular}
\end{table}

Die Klassen sind vor der Auswertung festgelegt, nicht danach. Der Grund ist methodisch: Wer
die Kategorien erst nach dem Blick auf die Fehlfälle bildet, passt sie unmerklich dem
Ergebnis an, und die Fehleranalyse verliert ihre Funktion als Kontrolle. Jeder Fall bekommt
eine Hauptursache, weil sonst ein einziger Fehler mehrfach gezählt würde; Nebenursachen
werden im Protokoll vermerkt.

Die Abgrenzung zur Abweichungstypologie ist wichtig. Ein Extraktionsfehler ist dort der
erste Typ und führt zum Ausschluss des Falls; in der Fehlertaxonomie ist er eine Messung der
Verfahrensqualität. Dieselbe Beobachtung kann also an zwei Stellen auftauchen, aber mit
verschiedener Funktion: einmal als Grund, einen Fund zu verwerfen, einmal als Kennzahl
dafür, wie oft das nötig war. Die Häufigkeit je Klasse wird erst nach dem Lauf berichtet;
sie ist eine Messzahl und steht deshalb in Kapitel~\ref{sec:ergebnisse} und nicht hier.

\subsection{Code-Verzeichnis}
\label{app:code}

Der Bericht wird aus dem Verzeichnis \texttt{report/} gesetzt. Die Auswertungswerkzeuge
entstehen in den Arbeitspaketen P2 und P3; die Verzeichnisse für Korpus, Code und Tests sind
noch nicht angelegt. Tabelle~\ref{tab:dateien} verzeichnet deshalb den heutigen Bestand und
keinen geplanten.

\begin{table}[ht]
  \centering
  \caption{Dateien und Verzeichnisse des Berichts}
  \label{tab:dateien}
  \small
  \begin{tabular}{@{}>{\raggedright\arraybackslash}p{0.30\textwidth}>{\raggedright\arraybackslash}p{0.62\textwidth}@{}}
    \toprule
    Datei oder Verzeichnis & Rolle \\
    \midrule
    \texttt{build.sh} & baut \texttt{main.pdf}; Ziele \texttt{expose} und \texttt{clean} \\
    \texttt{main.tex} & Wurzel; bindet Vorspann, Kapitel, Anhang und Literatur ein \\
    \texttt{literatur.bib} & Literaturangaben; 32 Einträge, davon 13 im Bericht zitiert \\
    \texttt{preamble/} & Pakete, Satzspiegel und Kommandos (drei Dateien) \\
    \texttt{frontmatter/} & Titelseite, Zusammenfassung, Selbstständigkeitserklärung \\
    \texttt{chapters/} & acht Kapiteldateien, \texttt{01\_einleitung} bis \texttt{08\_fazit} \\
    \texttt{backmatter/} & Anhang A.1 bis A.5 \\
    \texttt{figures/}, \texttt{tables/} & Ablage für Abbildungen und Tabellen; derzeit leer \\
    \texttt{expose.tex} & Ein-Seiten-Exposé als eigenes Dokument mit derselben Präambel \\
    \texttt{README.md} & Bauanleitung und Hinweise zur Engine \\
    \texttt{.gitignore} & Build-Artefakte; \texttt{main.pdf} wird mitgeführt \\
    \bottomrule
  \end{tabular}
\end{table}

Ein Code-Verzeichnis im engeren Sinn ist dieser Anhang noch nicht. Solange die Pipeline nicht
läuft, wäre eine Liste von Modulen eine Behauptung über Dateien, die es nicht gibt. Sobald
die Werkzeuge stehen, wird die Tabelle um ihre Einträge ergänzt; die Rollen der vorhandenen
Dateien bleiben dabei unverändert.
